\subsection{A.1. Objects, forms, and the three readings}

Throughout A.1--A.2, $k\subset\overline\Q\subset\C$ has its specified
embedding and $M$ belongs to Andr\'e's semisimple, graded category of pure
motivated motives with rational coefficients. Direct sums of pure weights are
allowed. The same arguments apply to a semisimple graded Tannakian category
with the indicated realizations, central weight cocharacter, and Hodge
property for its tensors. Reductivity is used below; these results do not
assert a corresponding dimension table for arbitrary mixed motives.
Andr\'e's Theorem~0.4, \S2.5 and \S4.6 supply this category, the Hodge
property, and the motivic group. His remark~(ii) in \S4.6 makes it necessary
to allow a nonconnected motivic group.
\footnote{Y. Andr\'e, \emph{Pour une th\'eorie inconditionnelle des motifs},
Publ. Math. IH\'ES 83 (1996), 5--49, Theorem 0.4, \S2.5, \S4.6 and
remark (ii); full text opened at
\url{https://www.numdam.org/item/10.1007/BF02698643.pdf}.}

\begin{result}{A.1.1}{The typed triangle and its homogeneous spaces}{Proved}
Put $V=T_{\B}$, $D=T_{\dR}$ and $H=\operatorname{gr}_F D$. On the weight
$n$ part, $\operatorname{gr}^p_F D$ is assigned bidegree $(p,n-p)$.
Thus $T_{\mathrm{Hdg}}$ here means the associated-graded de Rham tensor
fibre functor over $k$. The actual maps have types
\[
 \alpha:V_{\C}\longrightarrow D_{\C},\qquad
 \beta:D_{\C}\longrightarrow H_{\C},\qquad
 \gamma=(\beta\alpha)^{-1}:H_{\C}\longrightarrow V_{\C}.
\]
They are tensor isomorphisms; $\beta$ is the projection of each Hodge
summand onto its associated-graded summand. Write
\[
 G=\Aut^\otimes(V),\quad G'=\Aut^\otimes(D),\quad
 B'=\Aut^\otimes(D,F),\quad
 U'=\ker\bigl(B'\longrightarrow\Aut^\otimes(H)\bigr).
\]
The normalized splitting space
\[
 Q=\{s:D\xrightarrow{\sim}H\text{ tensor and filtered}:
                        \operatorname{gr}(s)=\id_H\}
\]
is a $U'$-torsor. Geometrically the filtration and bigrading spaces are
$G/B$ and $G/L$, where $B$ and $L$ are the transported parabolic and Levi
stabilizers. With $g=\dim G$, $b=\dim B$, $l=\dim L$, $u=\dim U$, one has
\[
                 g=l+2u,\qquad b=l+u,\qquad g-b=u.
\]
The morphism $\BigSpace\to\Fil$ is a torsor under the relative unipotent
group over $\Fil$ whose fibre at $F$ is $R_u(\operatorname{Stab}(F))$.
Every fibre has dimension $u$.
\end{result}

\begin{proof}
Strictness and tensor compatibility of the pure Hodge filtration make
$\operatorname{gr}_F$ an exact tensor fibre functor. The pure Hodge
decomposition splits this filtration, giving $\beta$ and the displayed
triangle. Over an algebraic closure choose a splitting cocharacter $\mu$.
Then $B'=P(\mu)$, $L'=Z(\mu)$, and $U'=R_u(P(\mu))$. On the Lie algebra
of a reductive group the nonzero $\mu$-weights occur in opposite pairs;
the zero-weight subspace is $\Lie L'$. This proves the dimensions.
For a nonconnected group the same calculation applies to its identity
component; the finite components do not change dimensions.

Two normalized splittings $s_1,s_2$ satisfy
$s_2^{-1}s_1\in U'$, and every element of $U'$ occurs this way. This
proves freeness and transitivity. Existence after algebraic closure is
the splitting of a parabolic by a Levi subgroup, so $Q$ is a torsor.
Its points correspond exactly to Levi subgroups: the stabilizer of the
grading $s^{-1}H$ is a Levi subgroup, every Levi occurs by unipotent
conjugation, and its stabilizer in $U'$ is trivial since the weights of
$\mu$ on $\Lie U'$ are strictly positive. Characteristic-zero unipotent
torsors are trivial over $k$, but no preferred $k$-point is chosen.

The orbit--stabilizer descriptions give $G/B$ and $G/L$. The fibre of
$G/L\to G/B$ is $B/L$, on which the unipotent radical of $B$ acts simply
transitively. Conjugating this group as the flag varies gives the
relative group just stated. This constructs the torsor action; it does
not require a single constant unipotent subgroup to act on all fibres.
\end{proof}

Here are the forms and the quotient maps precisely. The moduli functors
of tensor filtrations and tensor gradings give algebraic parameter
spaces. A fixed geometric cocharacter type has its field of definition,
the number field fixing its Galois orbit. Over that field take the
corresponding component; over $\Q$ take the finite union of conjugate
types. A rational stabilizer $B\subset G$ need not have been selected.
For the actual type over $k$, the torsor
$P=\Isom^\otimes(V_k,D)$ gives the filtration form as the quotient by
the left $B'$-action, with map
\[
                       q(a)=a^{-1}F_D.
\]
Similarly $P''=\Isom^\otimes(H,V_k)$ maps to the grading space by
transporting the fixed grading on $H$; call this map $r$. Thus
$\phi=q(\alpha)$, $\psi=r(\gamma)$ and $\psi$ maps to $\phi$.
Coordinates of these points in the indicated $\Q$-forms define
$\Q(\phi)$ and $\Q(\psi)$; algebraic extension of the type field
does not affect their transcendence degrees.

\begin{result}{A.1.2}{The $H'$ reading, the margin arrow, and the denominator}{Proved or refuted as specified}
The two descriptions of $Q$ in P11 [I] are related by a precise
normalization. If $H'$ is the inverse image in $B'$ of a nontrivial
central one-dimensional torus $T_0=\im(j)$ in $B'/U'$, the enlarged
splitting space $Q_H$ allowing $\operatorname{gr}(s)\in T_0$ is an
$H'$-torsor. There is a morphism
\[
 Q_H\longrightarrow T_0,\qquad s\longmapsto\operatorname{gr}(s),
 \qquad Q=\text{fibre over }1.
\]
Consequently $\dim Q_H=u+1$ and $\dim Q=u$. After a splitting,
$Q_H\simeq Q\times T_0$. If $j$ is an embedding, $T_0$ is the displayed
$\Gm$; with a finite kernel the correct quotient is its image torus.
The actual canonical $\beta$ lies in the normalized fibre $Q$.

The natural parameter-space arrow is $\Phi\to F$, forgetting the
grading. It induces an isomorphism on the discrete spaces of types;
the inverse on types can have direction $F/G\to\Phi/G$. A
$G$-equivariant section of the full parameter-space arrow exists only
when $u=0$. For $G=\GL_2$ and a line filtration, even an arbitrary
regular section is impossible.

The lowercase-$k$ reading in P17 [I] gives the field extensions in the
table below. The uppercase-$K$ reading of $\trdeg k(\alpha)/K$ or
$\trdeg k(\beta)/K$ is undefined unless the numerator equals $K$,
because $K=k(\alpha,\beta)$ contains both numerator fields; in that
exceptional case the degree is zero.
\end{result}

\begin{proof}
The quotient of two enlarged splittings belongs exactly to $H'$.
Taking associated graded gives the map to $T_0$, whose identity fibre
is the kernel torsor under $U'$. This proves the stated relation of
the two readings and the dimensions. The contracted-product reading
$P'=Q_H\times^{H'}G'$ consequently concerns the enlarged reduction;
its $u$-dimensional normalized fibre is the splitting space used below.

A $G$-equivariant section $G/B\to G/L$ would give a $B$-fixed point
of $B/L$. Its subgroup $U$ acts simply transitively there, so a fixed
point forces $U=1$. For $\GL_2$, the map is
$(\mathbf P^1\times\mathbf P^1)\setminus\Delta\to\mathbf P^1$.
A section would give a morphism $f:\mathbf P^1\to\mathbf P^1$ with
no fixed point. For a nonconstant $f=[P:Q]$, the equation
$XQ(X,Y)-YP(X,Y)=0$ is homogeneous of positive degree and has a
projective zero over an algebraic closure; if it vanishes identically,
every point is fixed. A constant map fixes its value. This is a
counterexample to the literal reversed arrow on parameter spaces.
The field assertion follows directly from the inclusions defining $K$.
\end{proof}

\begin{result}{A.1.3}{The Tate-period membership with the stated direction of $\alpha$}{Proved}
If $\Q(-n)$ belongs to the tensor category generated by $M$, then
$(2\pi i)^n\in k(\alpha)$.
\end{result}
\begin{proof}
For $\Q(-1)=H^2(\mathbf P^1)$, use the positive Betti generator and the
algebraic de Rham generator normalized by $d\log$. In the direction
$\alpha:T_{\B}\to T_{\dR}$ the coefficient is $(2\pi i)^{-1}$, hence
on $\Q(-n)$ it is $(2\pi i)^{-n}$. An algebraic projector and the
inclusion and projection of a tensor summand express its comparison
matrix rationally in that of $M$; inverses of matrices are rational
functions as well. Thus $(2\pi i)^{-n}\in k(\alpha)$, and inversion
proves the claimed membership. The normalization is also checked by
the explicit cycle-class computation in F.1.
\end{proof}

\subsection{A.2. The joint period torsor and the corrected dimension table}

\begin{result}{A.2.1}{Algebraic extension and the conjugate comparison}{Proved}
Passing from an algebraic $k\subset\C$ to $\Qb$ preserves the
identity component of the motivated group on the given Betti
realization, and hence $g,b,l,u$. It preserves every absolute and
relative transcendence degree used here. Over the compositum
$k\,c(k)$, and then over $\Qb$, the object
\[
                             N=M\oplus c^*M
\]
is defined. Conjugation gives a rational tensor map $j_{\B}:V\to V_c$
and a $c$-semilinear map $j_{\dR}:D\to D_c$. Their comparison identity is
\[
                   \alpha_c j_{\B}=j_{\dR}\alpha\kappa_V,
                   \tag{A.2.1a}
\]
where $\kappa_V$ is coordinate conjugation on $V_{\C}$. Both sides of
this equation are $\C$-linear. In bases carried by $j_{\B}$ and
$j_{\dR}$, the matrix of $\alpha_c$ is $\overline\alpha$.
\end{result}

\begin{proof}
A motivated tensor over an algebraic extension is constructed from
finitely many varieties, cycles and coefficients, hence over a finite
extension. Enlarge that extension to a finite Galois one. On the
fixed tensor construction in $M$, its conjugates form a finite set
of rational Betti tensors. The symmetric tensors in this set descend:
the algebraic cycles occurring in the motivated construction descend
after taking the finite traces and products giving those symmetric
tensors. The group over $k$ therefore preserves the polynomial with
these tensors as factors and permutes its finite set of factors.
Its identity component fixes each factor. It consequently fixes
every motivated tensor available after algebraic extension. The
group after extension is also a subgroup of the original group, so
their identity components coincide. The Lie-algebra Hodge-weight
decompositions give the remaining dimensions. This argument does not
assume geometric connectedness.

For any coordinate tuple $z$, $\Qb(z)/k(z)$ is algebraic. Invariance
of transcendence degree under algebraic extensions, and the same
fact in each field tower, proves the numerical assertion. Finally
conjugation of complex varieties gives $j_{\B}$, including its usual
Tate sign; conjugation of algebraic differential forms gives
$j_{\dR}$. Integrating conjugate forms over conjugate cycles gives
(A.2.1a). The two semilinear maps on its right compose to a linear
map, which checks its type.
\end{proof}

Let $T=P_N$ be the motivated comparison torsor, with projection
morphisms $a:T\to P_M$ and $d:T\to P_{c^*M}$. Identify the second
Betti group with $G$ through $j_{\B}$. The joint Betti group is a
closed subgroup
\[
                 J\subset G\times G
\]
whose projections are surjective. Set $N_2=\ker(\operatorname{pr}_1)$,
viewed in the second $G$; it is normal there. Define on $T$
\[
 F_1(t)=a(t)^{-1}F_D,\qquad
 F_2(t)=j_{\B}^{-1}d(t)^{-1}F_{D_c}.
 \tag{A.2.1b}
\]
Let $T^\circ$ be the open locus of pairs opposed in every motivic
weight. At the actual comparison $t_0=(\alpha,\alpha_c)$, equation
(A.2.1a) makes $F_2=\overline{F_1}$, so purity proves
$t_0\in T^\circ$. In particular this open locus is nonempty.
Its existence has been proved from the actual comparison; an
algebraic comparison point was not presumed to have this property.
Once nonemptiness is known, a nonempty open variety over $\Qb$ has
$\Qb$-points and can be used for algebraic trivializations.

On $T^\circ$, the intersections
$F_1^p\cap F_2^{n-p}$ define a tensor bigrading. Transport by $a(t)$
and projection to $\operatorname{gr}_F D$ give $\beta(t)\in Q$.
Define the remaining maps by
\[
 \gamma(t)=(\beta(t)a(t))^{-1},\quad
 \phi(t)=F_1(t),\quad
 \psi(t)=\bigl(F_1^p(t)\cap F_2^{n-p}(t)\bigr)_{p,n}.
 \tag{A.2.1c}
\]
These are regular: intersections of the indicated constant ranks and
their projections are expressed using invertible minors on the
opposed open set. Their targets are respectively $P_M,Q,P'',\Fil$
and $\BigSpace$, and they recover the actual five points at $t_0$.
In particular the splitting construction proves, over the compositum,
\[
             k(\alpha,\beta)\subset k(\alpha)\,c(k(\alpha)).
 \tag{A.2.1d}
\]

\begin{result}{A.2.2}{All six algebraic image dimensions}{Proved}
Use $t_0$ only to identify geometric groups and flags over $\C$.
Let $\mu$ be the Hodge cocharacter of $M$, $w$ the central weight
cocharacter, and write
\[
 B=P_G(\mu),\quad B^-=P_G(w\mu^{-1}),\quad
 L=Z_G(\mu),\quad U=R_u(B).
\]
The joint Hodge cocharacter is
$\lambda=(\mu,w\mu^{-1})\in J_{\C}$. Put
\[
 n=\dim N_2,\quad n_B=\dim(N_2\cap B^-),\qquad v=n-n_B.
 \tag{A.2.2a}
\]
Then $0\leq v\leq u$. On the component containing $t_0$, the
dimensions of the closures of the images in (A.2.1c) are
\begin{center}
\begin{tabular}{c c c c c c}
\toprule
$\alpha$&$\beta$&$\gamma$&$\phi$&$\psi$&$(\alpha,\beta)$\rule{0pt}{2.4ex}\\
\midrule
$g$&$u$&$b+v$&$u$&$u+v$&$g+v$\\
\bottomrule
\end{tabular}
\end{center}
More precisely, under the algebraic coordinate change
\[
 P\times Q\xrightarrow{\sim}Q\times P'',\qquad
 (a,s)\longmapsto(s,(sa)^{-1}),
 \tag{A.2.2b}
\]
the closure of the joint image is $Q\times\Gamma$, where $\Gamma$
is the closure of the $\gamma$-image and $\dim\Gamma=b+v$.
\end{result}

\begin{proof}
The Hodge cocharacter belongs to the joint group because motivated
tensors are Hodge tensors. Its parabolic is
\[
                   B_J=P_J(\lambda)=J\cap(B\times B^-).
\]
The first projection $B_J^0\to B^0$ is surjective. Indeed the
surjective map $\Lie J\to\Lie G$ respects cocharacter weights and
is surjective on each weight, hence on the nonnegative-weight
subspaces defining the parabolics. The image of the connected
parabolic is a closed subgroup of full Lie algebra in $B^0$, so it
is $B^0$. Its kernel has dimension $n_B$. The same reasoning works
with the relevant full stabilizers in the nonconnected groups.
Also $\dim J=g+n$.

Projection to $a$ has image a torsor component of dimension $g$.
Its filtration quotient has dimension $g-b=u$. The pair of flags
has stabilizer $B_J$, so the bigrading image, obtained by restricting
to the opposed open set, has dimension
\[
 (g+n)-(b+n_B)=u+v.
\]
The number $v$ counts the strictly negative Hodge weights of
$\Lie N_2$, so $0\leq v\leq u$. Equivalently it is the dimension
of the $N_2$-orbit in $G/B^-$. Normality of $N_2$ makes that orbit
dimension the same at every flag in the $G$-orbit.

For the remaining fibres, make the following subgroups first in the
de Rham form of $J$, using its algebraically defined filtration
stabilizers, and identify them geometrically with
\[
 S_\gamma=J\cap(U\times B^-),\qquad
 S_K=J\cap(1\times B^-).
\]
Here is an exact fibre calculation. Write $t'=(xa,yd)$ using the
simply transitive action of the de Rham joint group. If
$x\in U$, $y\in B^-$, the second Betti flag is unchanged; the
first is unchanged because $x$ preserves $F_D$. The opposite de Rham
flag is transported by $x$, so its normalized splitter is
$\beta'=\beta x^{-1}$. Consequently $\gamma'=\gamma$.
Conversely $\gamma'=\gamma$ gives
$\beta'a'=\beta a$, and hence
$x=a'a^{-1}=\beta'^{-1}\beta\in U$. The common $\gamma$
transports the fixed grading of $H$ to the same Betti bigrading;
in particular it determines the same second flag. Equation
(A.2.1b) then gives $y\in B^-$. Thus each $\gamma$-fibre is exactly
an $S_\gamma$-coset. Requiring also $a'=a$ gives precisely an
$S_K$-coset for the $(a,\beta)$-fibre.

The projection $S_\gamma\to U$ is surjective, since $U$ is
connected and $B_J^0\to B^0$ is surjective. Its kernel is $S_K$.
Therefore
\[
 \dim S_K=n_B,\qquad \dim S_\gamma=u+n_B,
\]
which yields $\dim\im(a,\beta)=g+v$ and
$\dim\im\gamma=g-u+v=b+v$. On a fixed $\gamma$-fibre, the
projection to normalized splittings is the full simply transitive
$U$-action. It ranges through every point of $Q$. Applying
(A.2.2b) thus gives the product of $Q$ with the $\gamma$-image;
taking closures proves the product assertion and
$\dim\im\beta=u$. This proves all dimensions without a
transcendence assumption on the actual splitting. Finite components
alter none of the calculations.
\end{proof}

\begin{result}{A.2.3}{When the image fills $P\times Q$}{Proved}
The image fills its component of $P\times Q$ exactly when $v=u$.
Equivalently the projected Hodge cocharacter is central on the
identity component of the common quotient $G/N_2$, or the Hodge
parabolic of that quotient has full dimension.
\end{result}
\begin{proof}
The ambient component has dimension $g+u$, whereas its irreducible
image closure has dimension $g+v$. Equality of these dimensions is
equivalent to equality of the closures. On the quotient Lie algebra,
the dimension of the positive Hodge-weight part is $u-v$.
Reductivity pairs its positive and negative weights, so this
vanishes exactly when all its weights are zero. This is centrality
of the projected cocharacter on the identity component. One may
replace $N_2$ by $N_2^0$ for this dimension criterion; finite quotients
do not change it. In particular a shared Tate or other central torus
does not obstruct $v=u$.
\end{proof}

\begin{result}{A.2.4}{Actual-point bounds and the requested period-conjecture implications}{Proved}
The actual degrees have the unconditional upper bounds in the middle
column below. The last column gives their exact values under
$\mathrm{PG}_{\mathrm{mot}}(N)$, the assertion that the comparison
point of the named motive $N=M\oplus c^*M$ is generic in its
motivated comparison-torsor component.
\begin{center}
\begin{tabular}{lcc}
\toprule
Degree&Unconditional upper bound&Under $\mathrm{PG}_{\mathrm{mot}}(N)$\\
\midrule
$\trdeg_k k(\alpha)$&$g$&$g$\\
$\trdeg_k k(\beta)$&$u$&$u$\\
$\trdeg_k k(\gamma)$&$b+v$&$b+v$\\
$\trdeg_{\Q}\Q(\phi)$&$u$&$u$\\
$\trdeg_{\Q}\Q(\psi)$&$u+v$&$u+v$\\
$\trdeg_k K$&$g+v$&$g+v$\\
$\trdeg_{k(\alpha)}K$&$v$&$v$\\
$\trdeg_{k(\beta)}K$&$b+v$&$b+v$\\
$\trdeg_{k(\gamma)}K$&$u$&$u$\\
\bottomrule
\end{tabular}
\end{center}
This assertion of genericity is a consequence of the Grothendieck
period conjecture in the torsor formulation of Bost--Charles,
Conjecture~2.12, for a smooth projective generator of $N$.
The passage uses their Definitions~2.4 and~2.9, Corollary~2.11 and
Propositions~2.13--2.14: the comparison closure lies in the motivated
torsor, which lies in the algebraic-cycle torsor. Equality with the
last makes the comparison generic in the intermediate torsor; the
projection to the tensor subcategory generated by $N$ is dominant.
\footnote{J.-B. Bost and F. Charles, \emph{Some remarks concerning the
Grothendieck period conjecture}, \S2, exact locators just listed; full
author text opened at
\url{https://www.math.ens.psl.eu/~charles/crelle-periods.pdf}.}
\end{result}

\begin{proof}
Every image closure in A.2.2 is defined over an algebraic field.
The transcendence degree of a point's coordinates is the dimension
of its Zariski closure, so the six absolute bounds follow. For
relative bounds it would be insufficient to use only a generic
fibre. The product $Q\times\Gamma$ proves for every point that
adjoining $\gamma$ to $\beta$ adds at most $b+v$, and adjoining
$\beta$ to $\gamma$ adds at most $u$. The inverse of (A.2.2b) shows
that these joint fields equal $K$.

Fixing an arbitrary first comparison $a$, the second comparison
varies under $N_2$. Its flag image is an open subset of an
$N_2$-orbit in $G/B^-$. Every such orbit has dimension $v$, by
normality. Thus the fibre over the actual coordinate field
$k(\alpha)$ has dimension $v$, proving
$\trdeg_{k(\alpha)}K\leq v$ even when $\alpha$ is nongeneric.
Equivalently, $T/S_K$ is a homogeneous bundle over $P$ with these
fibres, and $T^\circ$ imposes an open condition.

A rational map defined at a generic comparison point sends it to
the generic point of its image: the induced map on image function
fields is injective. Apply this to all six maps to obtain the exact
absolute values under $\mathrm{PG}_{\mathrm{mot}}(N)$. Subtraction
in the three field towers gives the last three values. A.2.1
transfers all these degrees back to the original $k$.
\end{proof}

The comparison of this table with the list on 18:11 is exhaustive
at the level of the proved algebraic bounds and the requested
conjectural implications. The original $\alpha$ and $\phi$ values
already follow from $\mathrm{PG}_{\mathrm{mot}}(M)$ alone.
The original $\beta$ and $K/k(\gamma)$ predictions agree with the
table for every $v$. The five original values for $\gamma$, $\psi$,
$K$, $K/k(\alpha)$ and $K/k(\beta)$ require $v=u$; for $v<u$
they are each false unconditionally, since their proposed values
exceed the corresponding unconditional bound. When $v=u$, all
original values follow from $\mathrm{PG}_{\mathrm{mot}}(N)$.
This does not prove actual maximality for every motive.

There are further consequences from $\mathrm{PG}_{\mathrm{mot}}(M)$
alone. Since $\trdeg_k K\geq g$, the uniform relative bounds imply
\[
 \trdeg_k k(\beta)\geq u-v,\qquad
 \trdeg_k k(\gamma)\geq b.
\]
For $v=0$ this determines every entry of the corrected table.
For other $v$ it states the precise lower bounds obtained here;
no claim is made that other arithmetic results cannot strengthen
them.

\begin{result}{A.2.5}{Real models and Lead A}{Proved}
Suppose $M$ is obtained from a motive over an algebraic
$k_0\subset\R$. Let $f_\infty$ be its rational cohomological real
Frobenius. Then, as equations of flags,
\[
 \overline\phi=f_\infty\phi,\qquad
 F_D^{\mathrm{opp}}=\alpha f_\infty\alpha^{-1}F_D.
\]
Consequently
\[
                     \Q(\psi)=\Q(\phi),\qquad K=k(\alpha).
 \tag{A.2.5a}
\]
Over $\Qb$ one has $v=0$, and the corrected conjectural tuple is
\[
 (\alpha,\beta,\gamma,\phi,\psi,K;K/\alpha,K/\beta,K/\gamma)
                 =(g,u,b,u,u,g;0,b,u),
 \tag{A.2.5b}
\]
where entries denote the degrees with the bases in A.2.4 and
$\mathrm{PG}_{\mathrm{mot}}(M)$ suffices for these values.
\end{result}

\begin{proof}
Conjugation exchanges the Hodge indices, giving the first flag
identity. The bigrading is recovered by intersections of that flag
with its $f_\infty$-image, hence is rational over $\Q(\phi)$;
the forgetful map gives the reverse field inclusion. Transport of
the same opposite flag through $\alpha$ gives the second identity.
Intersections and projections are rational operations, so $\beta$
and $\gamma$ are defined over $k(\alpha)$, proving (A.2.5a).
The involution need not lie in the smaller motivic group after a
nonreal base extension; it normalizes it by conjugating the tensor
conditions, which is sufficient for these equations.

Over $\Qb$, $M\simeq c^*M$; the tensor categories generated by
$M$ and by their direct sum coincide. The joint group is a graph,
so $N_2=1$ and $v=0$. Apply A.2.4 and the subsequent
$\mathrm{PG}_{\mathrm{mot}}(M)$ lower bounds.
\end{proof}

The same dimension conclusion $v=0$ holds when $M\simeq c^*M$
over $\Qb$ without a chosen real model. Defining an isomorphism may
require algebraic extension of $k$, so the unextended identity
$K=k(\alpha)$ must then be replaced by that identity after the
extension; the transcendence degrees are unchanged. A fixed such
isomorphism induces a rational Betti map carrying the conjugate
Hodge flag to the original one, which also yields
$\Q(\psi)=\Q(\phi)$. When no such isomorphism is available the
joint-torsor construction and the parameter $v$ still apply.

For $u=0$, $Q$ has dimension zero and its normalized splitting is
algebraic; the filtration and bigrading points are algebraic as
well. Then $K=k(\alpha)=k(\gamma)$ up to an algebraic extension,
and the only remaining absolute period degree is the ordinary
degree of the comparison point. Thus the degenerate case is
included without an assumption of a real model.

\begin{result}{A.2.6}{An explicit $\GL_2$ refutation and Lead H2}{Proved; remaining equivalence specified}
Let $k=\Q$, use its real embedding, and put
\[
                 E:y^2=x^3-x+1,\qquad M=h^1(E).
\]
Its discriminant is $-368$ and $j(E)=-6912/23$. A CM elliptic
$j$-invariant is an algebraic integer, so this curve is non-CM.
Thus $G=\GL_2$, $g=4$, $b=3$, $l=2$, $u=1$, and $v=0$.
For this completely specified motive, the five excessive original
values described above are false independently of any period
conjecture. In particular conjugate periods add zero new
transcendentals to $k(\alpha)$, whereas $u=1$.
\end{result}

\begin{proof}
For the group assertion, $\End_{mathrm{HS}}H^1(E,\Q)
=\End(E)\otimes\Q=\Q$: a rational map preserving the complex
structure becomes a complex-torus homomorphism after clearing
lattice denominators. The Mumford--Tate representation is
irreducible and contains scalar homotheties. A connected reductive
proper subgroup of $\GL_2$ with irreducible two-dimensional
representation would have to be a nonsplit torus, whose commutant
is a quadratic field; this contradicts the displayed endomorphism
algebra. Thus its group is $\GL_2$. Motivated tensors are Hodge
tensors, so the motivated group contains this group and hence is
also $\GL_2$. The arithmetic test for CM uses the integral Hilbert
class polynomial; a checked explicit locator is Sutherland's
MIT 18.783 Lecture 20 (20 April 2022), Theorem 20.12 and Corollary
20.13, printed p.~6.
\footnote{Opened author lecture text:
\url{https://math.mit.edu/classes/18.783/2022/LectureNotes20.pdf}.}

Choose a rational Betti eigenbasis for
$f_\infty=\operatorname{diag}(1,-1)$ and a de Rham basis with
$F^1=\C e_1$. Put
\[
 A=\begin{pmatrix}a&b\\c&d\end{pmatrix},\qquad \delta=ad-bc.
\]
At the actual comparison $c,d\ne0$: otherwise the Betti Hodge
line is an eigenline of $f_\infty$ and equals its conjugate,
contradicting opposedness. Exact multiplication gives
\[
 Af_\infty A^{-1}=\frac1\delta
  \begin{pmatrix}ad+bc&-2ab\\2cd&-ad-bc\end{pmatrix},
 \qquad
 \beta=\begin{pmatrix}1&-\dfrac{ad+bc}{2cd}\\0&1\end{pmatrix},
\]
\[
 \beta A=\begin{pmatrix}\delta/(2d)&-\delta/(2c)\\c&d\end{pmatrix},
 \qquad
 \gamma=\begin{pmatrix}d/\delta&1/(2c)\\-c/\delta&1/(2d)\end{pmatrix}.
\]
In particular
\[
               \gamma_{11}\gamma_{22}+\gamma_{12}\gamma_{21}=0.
\]
The image has dimension three, also visible by recovering $c,d$
and $\delta$ rationally from $\gamma$; they are independent
coordinates on the displayed open image. This proves the
$\gamma$ bound directly as well as by A.2.2. The remaining
failures follow from $K=k(\alpha)$ and $\Q(\psi)=\Q(\phi)$
and the uniform bounds of A.2.4. Exact symbolic programs,
inputs and output are retained in the computation record; no
numerical approximation is used in this counterexample.
\end{proof}

Lead H2 interpreted as the implication from joint period
genericity to maximal splitting degree is proved by the dominant
$\beta$-map of A.2.2. Its proposed reason that the conjugate periods
always add exactly $u$ transcendentals is refuted by A.2.6. The
correct algebraic object is the image of $T^\circ\to P\times Q$,
of dimension $g+v$; it need not be the whole product or the whole
joint torsor. Dominance alone does not prove a converse from
genericity of one projection to genericity of the joint point.
We have not produced a motive for which those two arithmetic
genericity assertions have provably different truth values.
Consequently H2 interpreted as that literal arithmetic equivalence
remains an unresolved obligation in this incomplete record.

\subsection{A.3. Transfer to established period-conjecture cases}

\begin{result}{A.2.7}{Every established joint genericity case gives the corrected table}{Proved}
For every specific motive $M$ for which
$\mathrm{PG}_{\mathrm{mot}}(M\oplus c^*M)$ is established, all nine
equalities in A.2.4 hold unconditionally. This is a theorem about
every such case, since its proof is the application of the six
fixed rational maps, not a new transcendence hypothesis.
\end{result}
\begin{proof}
Insert the established genericity into A.2.4. Its argument proves
injectivity of the function-field maps and the relative tower
equalities for that named motive. Algebraic base extension is
already handled by A.2.1. Thus no further comparison conjecture
is used for any of the nine entries.
\end{proof}

For a CM elliptic curve, $g=2$, $u=0$, and the curve is isogenous
to its conjugate over $\Qb$. The Chudnovsky input checked in the
elliptic section consequently gives the tuple
$(2,0,2,0,0,2;0,2,0)$. The same transfer applies to a tensor
subcategory generated by this curve and Artin objects; passage to
a finite coefficient extension removes the latter and introduces
no transcendence. It is not an assertion of a period conjecture
for arbitrary products of unrelated CM elliptic curves.

For an arbitrary elliptic curve, the line and ordered pair of
conjugate lines give, in a rational homology basis,
$\Q(\phi)=\Q(\tau)$ and
$\Q(\psi)=\Q(\tau,\overline\tau)$. If the curve is isogenous
to its conjugate, the isogeny yields a fractional linear
$\GL_2(\Q)$ relation between $\tau$ and $\overline\tau$, so
these two fields coincide. The opened Schneider and elliptic
period sources, and the unresolved case when the conjugate curve
is not isogenous, are treated in the following elliptic section;
the dimension calculation here does not establish that remaining
two-variable algebraic independence.
